同一份代码，同一个步长，只换了随机种子，两次训练分别停在损失 \(2.31\) 和 \(2.05\)。两次的梯度范数都降到了 \(10^{-4}\)。

哪一次算是收敛了？两次都停住了，可它们停在不同的地方，而你手上没有任何东西能判断 \(2.05\) 离最好还有多远——也许它就是最优，也许还差得远。\textbf{局部最优是站在原地能确认的事，全局最优是关于整个可行域的断言}，一般问题里前者不蕴含后者。

凸性做的正是这一件事：它让两者合而为一。而支撑这个结论的，是一条形状朴素的不等式——

\[
f(y)\;\ge\;f(x)+\nabla f(x)^{\trans}(y-x).
\]

在一点算出的梯度，给出了对\emph{所有}其他点的下界。局部信息第一次有了全局效力。这一部分的其余内容，几乎都是这条不等式的延伸：最优性怎么认证、下界怎么抬高、算法每一步在改进什么、以及这一切在非凸地形上还剩多少。

\begin{center}
\begin{tikzpicture}[x=1cm,y=1cm,font=\small,>=stealth]
  \draw[black!45,fill=blue!9] (0,1.4) rectangle (3.5,2.9);
  \node[align=center] at (1.75,2.15) {切平面\\是全局下界};
  \draw[black!45,fill=green!8] (5.0,3.5) rectangle (13.2,4.4);
  \node[align=left,anchor=west] at (5.2,3.95) {梯度为零 \(\Rightarrow\) 全局最优（不必再搜）};
  \draw[black!45,fill=green!8] (5.0,2.35) rectangle (13.2,3.25);
  \node[align=left,anchor=west] at (5.2,2.8) {对偶：可行乘子给出可验证的下界};
  \draw[black!45,fill=green!8] (5.0,1.2) rectangle (13.2,2.1);
  \node[align=left,anchor=west] at (5.2,1.65) {算法：每步用局部模型换取确定的进展};
  \draw[black!45,fill=red!7] (5.0,0.05) rectangle (13.2,0.95);
  \node[align=left,anchor=west] at (5.2,0.5) {非凸：以上三条全部降级，只剩驻点};
  \draw[->,black!55] (3.55,2.3) -- (4.95,3.95);
  \draw[->,black!55] (3.55,2.2) -- (4.95,2.8);
  \draw[->,black!55] (3.55,2.0) -- (4.95,1.65);
  \draw[->,dashed,black!45] (3.55,1.75) -- (4.95,0.5);
  \node[align=center,font=\footnotesize,black!65] at (1.75,0.5) {这条不等式\\一旦失效，\\右下角就是\\你的处境};
\end{tikzpicture}

\emph{九章分成两半：前七章讲这条不等式成立时能拿到什么，最后两章讲它不成立时还剩什么。深度学习的日常在下面那一格。}
\end{center}

\subsection{五条贯穿始终的判断}

\textbf{难的是翻译，不是求解。}把一段现实诉求写成决策变量、目标函数、约束三件套，这一步做错了，后面解得再准也是错的。目标写歪，最优解会精确地钻那个空子——只优化点击总数就催生夸张标题，只考核平均处理时长就催生提前结单。这与奖励塑形的失败是同一形状。

\textbf{凸性的操作性读法是``可以安全混合''。}任取两个可行方案，任意配比的折中仍然可行；任取两点，弦落在函数图像上方。这两句几何描述，比任何定义都更适合用来判断一个新问题凸不凸。判断的实用办法也不是求 Hessian，而是拆解构造过程：非负加权和、逐点上确界、仿射复合各自保凸，常见损失都能归到某条规则上。

\textbf{最优性需要证书，而对偶提供它。}原问题的可行解给上界，对偶的可行乘子给下界，两者一夹就知道离最优还有多远——这个判断不需要知道真正的最优值。求解器日志里的 gap 就是它。互补松弛进一步说明哪些约束真正在起作用，这解释了支持向量机为什么只依赖少数样本。

\textbf{算法的强弱取决于它调用哪种信息。}只用梯度，步长受曲率限制、速度受条件数限制；用上 Hessian，方向被校正且对坐标缩放不变，代价是每步要解一个线性系统；不可微时退到次梯度，收敛降级但仍能前进；结构可分时用近端算子把光滑与非光滑部分分开处理。接口越强，单步越贵，速率越快——这是一条贯穿全章的兑换关系。

\textbf{非凸不等于难优化，但保证确实降级了。}梯度法在非凸地形上只能保证梯度范数趋于零，既不保证全局最优、也不保证局部极小。高维里鞍点远多于局部极小；而某些非凸问题（满足 PL 条件的那些）照样有线性速率。更要紧的是把三类误差分开：训练损失降不下去是优化问题，训练低而测试高是统计问题，两者混淆会导致该加数据的时候去调优化器。

\subsection{九章怎么安排}

\hyperref[ch:073]{``凸函数：局部信息为何能控制全局''}是全部分的枢纽，其中``切平面为何是全局下界''一篇必须读透——后面所有章节都在用它。\hyperref[ch:072]{``凸集合：可行方案能否安全混合''}提供可行域一侧的对应结论，两章合起来是地基。

\hyperref[ch:075]{``对偶：为最优性建立下界证书''}是回报最高的一章。做机器学习的读者若只读两章，选它和凸函数那章：正则化系数与约束预算的等价、支持向量的稀疏性、求解器报告的 gap，都从这里出来。

\hyperref[ch:076]{``优化算法：怎样把局部模型变成全局进展''}解释训练中每天遇到的现象——学习率为何有上界、条件数为何决定速度、Adam 在模仿什么。\hyperref[ch:078]{``非凸与随机优化：能下降不等于找到全局解''}紧接其后，是深度学习优化的数学解释，两章建议连读。

\hyperref[ch:071]{``从现实选择到优化问题''}篇幅短但值得先看，它讲的是建模而非数学。\hyperref[ch:074]{``识别凸优化问题''}与\hyperref[ch:077]{``建模案例：凸性怎样改变真实问题''}偏应用，可按需取用；后者里 L1 稀疏那一篇解释了为什么会出现精确的零，值得单独一读。

\hyperref[ch:079]{``计算复杂度与近似算法：当最优在计算上不可及''}放在最后，回答一个立项阶段的问题：这个需求能不能要求``最优''。NP-hard 不意味着做不了，意味着放弃这个承诺，改求近似、启发式，或利用真实实例往往不是最坏情形。

读这一部分时值得反复问：我现在手上有的是一个\emph{解}，还是一个关于这个解有多好的\emph{证明}？凸优化的全部价值，在于它让第二样东西成为可能。
